% Revision 2, formatted from the user-supplied arXiv source archive.
\documentclass[11pt]{article}
\usepackage[utf8]{inputenc}
\usepackage{times}
\renewcommand{\ttdefault}{cmtt}
\usepackage[margin=1in]{geometry}
\usepackage{microtype}
\usepackage{amsmath,amsthm,amssymb}
\usepackage{graphicx}
\usepackage{booktabs}
\usepackage{array}
\usepackage{tabularx}
\usepackage{longtable}
\usepackage{xcolor}
\usepackage{caption}
\usepackage{flafter}
\usepackage{float}
\usepackage{placeins}
\renewcommand{\topfraction}{0.9}
\renewcommand{\bottomfraction}{0.7}
\renewcommand{\textfraction}{0.08}
\renewcommand{\floatpagefraction}{0.8}
\makeatletter
\setlength{\@fptop}{0pt}
\setlength{\@fpsep}{12pt plus 1fil}
\setlength{\@fpbot}{0pt plus 1fil}
\makeatother
\setlength{\textfloatsep}{14pt plus 2pt minus 4pt}
\setlength{\floatsep}{12pt plus 2pt minus 2pt}
\PassOptionsToPackage{hyphens}{url}
\usepackage{xurl}
\usepackage[colorlinks=true,linkcolor=black,citecolor=blue!50!black,urlcolor=blue!50!black]{hyperref}
\hypersetup{pdftitle={Auditing Privacy Claims in Split Language-Model Training},pdfauthor={Georgios Politis and Evangelos Pappas},pdfsubject={Retrospective audit of split-model privacy experiments}}
\newtheorem{proposition}{Proposition}
\captionsetup{font=small,labelfont=bf,labelsep=period}
\expandafter\newcommand\csname artifactnameE1\endcsname{decoy-gradient records}
\expandafter\newcommand\csname artifactnameB1\endcsname{boundary-audit records}
\expandafter\newcommand\csname artifactnameR1\endcsname{exploration records}
\expandafter\newcommand\csname artifactnameP4\endcsname{pilot records}
\expandafter\newcommand\csname artifactnameE8\endcsname{wrapper records}
\expandafter\newcommand\csname artifactnameER\endcsname{deployment records}
\expandafter\newcommand\csname artifactnameMI\endcsname{MI-analysis records}
\DeclareRobustCommand{\artifact}[1]{\hyperlink{artifact-#1}{\csname artifactname#1\endcsname}}
\newcommand{\pp}{\,\mathrm{pp}}
% Generated by analysis/recompute.py; retrospective archived-evidence counts.
\newcommand{\RoneCells}{19}
\newcommand{\RoneQualifiedBreaches}{13}
\newcommand{\RoneUtilityFailures}{5}
\newcommand{\RoneGainFailures}{7}
\newcommand{\RoneUndetermined}{6}
\newcommand{\PfourCells}{7}
\newcommand{\PfourUtilityFailures}{2}
\newcommand{\PfourPrivacyFailures}{4}
\newcommand{\PfourPrivacyUndetermined}{3}
\newcommand{\PinnedInputs}{122}
\newcommand{\ConsistencyChecks}{1,522}

\title{\textbf{Auditing Privacy Claims in Split Language-Model Training}}
\author{Georgios Politis\\ {\small\texttt{geo.politis@gmail.com}} \and
Evangelos Pappas\\ {\small\texttt{evanz.pappas@gmail.com}}}
\date{}
\begin{document}
\maketitle
\begin{abstract}
Privacy evaluations can fail for several reasons. An attack may miss a
visible message, score the wrong target or lack the controls needed to
support its verdict. A defence can also fail a sound test. We distinguish
these outcomes in a retrospective audit of four split-language-model
campaigns from one research programme, with two further analyses of
security scope and mutual information. The audit separates measured
breaches, qualified non-detection, incomplete evidence, invalid
measurements and claims made outside the tested threat model. In an
exploration of 19 configurations, only four of 24 required instruments
qualified, yet 13 configurations had qualified privacy breaches.
The six configurations with undetermined privacy outcomes independently
failed utility or gain requirements. A subsequent seven-cell pilot
produced no eligible defence, with two utility failures and four failed
versus three undetermined composite privacy assessments. These results
retain valid evidence against the tested defences. The scope analyses
show why shared-key wire obfuscation does not protect against a
compromised restoring host, and why a near-zero mutual-information lower
bound cannot certify a small leakage budget. We provide rerunnable
secondary analysis, hashed source records and each campaign's uncertainty
conventions. The resulting cross-case audit identifies the conclusions
supported by the recorded experiments. Its scope remains retrospective,
with no new defence or independently validated general evaluator.
\end{abstract}

\section{Introduction}
\label{sec:intro}

A split language model runs some layers on a client and others on a
remote provider. The client keeps its text locally, but the activations
and gradients exchanged during training can still reveal information
about it. Recovery scores need a clear context. Their meaning depends on
the messages the attacker sees, the model weights and secrets it knows,
the target it tries to recover and the controls used to assess its
performance. An evaluation can therefore report low leakage while testing
the wrong target or overlooking an available channel. Conversely, a
qualified attack that detects substantial leakage can provide sound
evidence against a defence.

Existing studies illustrate several parts of this problem. An omitted
returned-gradient channel can overturn a passing forward-only check
\cite{paper1}. Access to gradients also affects token recovery and
exact-document recovery by different amounts \cite{gradients2026}.
The scoring unit changes the result. LLM-Attack Suite records attacker
observations, controls, saved predictions and untested attacks
\cite{suite2026}, while work on image inversion shows how the evaluator
and attack influence measured leakage \cite{mia2026}. Together, these
findings motivate a closer account of which conclusions follow from
an experiment and which require evidence it never collected.

We examine four experimental campaigns within one research programme.
The cases cover an omitted gradient channel, a reconstruction test that
scored next-token prediction, incomplete instrument qualification and a
calibrated defence-redesign pilot. Their outcomes differ. Two further
analyses examine shared-key wire protection under a stronger observer
and the use of a near-zero mutual-information lower bound as a privacy
ceiling. Each case retains its own target, repetition unit and recovery
measure. Shared data and research decisions prevent treating them as
independent replications of one hypothesis or pooling their scores into
a common effect size.

The audit follows each privacy claim from the attacker's observations to
the recorded decision. We distinguish valid negative evidence from
non-detection, incomplete qualification, a broken measurement and a
change of threat model. The distinction has practical consequences.
Our executable analysis reconstructs the exploration's cell and seed
results and the defence pilot's saved losses and probe scores, preserving
their frozen rules and composite verdicts. Archived key, identity and
runtime records then delimit the broader security interpretations.
This reconciliation makes the existing evidence easier to inspect and
reproduce, without presenting the controls or attack infrastructure as
new methods.

\section{Scope and retrospective audit design}
\label{sec:design}

We audited claims from four campaigns in one split language model research
programme, then checked two broader claims about security scope. A claim
and its supporting records form the unit of analysis. We selected the
campaigns used in the original synthesis because they let us trace how
individual results became broader conclusions. The selection serves that
purpose, but provides no estimate of how often privacy evaluations fail
across the field. Each component result remains attributed to its source
campaign. The audit reconciles those records with the conclusions they
support.

The historical outcomes were available before we specified this audit.
We distinguish frozen exploration, completed experiments, proposed
confirmation and the present secondary analysis throughout. Our work
recomputes stored arithmetic and examines existing evidence without
retraining victim models, collecting captures, fitting stronger attackers
or filling gaps with inferred outcomes. Independent replication and
prospective confirmation would require their own studies. The
LLM-Attack Suite supplies relevant evaluation contracts and
controls~\cite{suite2026}. Here we use that context to examine claims
across campaigns, with no new validation of the suite itself.

\paragraph{Evidence required for a claim.}
For each claim, we record the protocol and cohort, the adversary's access
to secrets and remote state, and the channels it can observe. We also
identify the protected target and its alignment with observations, the
attack and training budgets, and the score, baseline and acceptance rule.
The qualification binding and the independent unit used to calculate
uncertainty complete the record. A qualification binding specifies the
target, channel, defence condition and implementation covered by positive
and null controls. A control passed elsewhere in the programme supplies
no qualification for a new binding. Likewise, recording a channel alone
does not evaluate attacks against it.

Matching generated records and currently admitted manifests take
precedence over narrative summaries. When a later decision supersedes an
earlier one, we retain that history together with the limitations attached
to the current status. This rule applies whichever decision is more
favourable. The source index links each retained quantity to its artifact
and distinguishes new arithmetic from results read through an existing
audit or manuscript. These links establish traceability within the supplied
package. Completeness of the original collection remains a separate
question.

\paragraph{Five outcomes for a specified claim.}
Different endpoints of the same configuration can have different
outcomes. We use the following terms consistently.
\begin{description}
  \item[Breach.] An admissible, qualified measurement exceeds the
    declared privacy rule. It establishes a breach for the specified
    observer, target, attack and finite protocol, without necessarily
    quantifying every disclosure.
  \item[Non-detection.] An admissible, qualified test stays within its
    rule for the tested attack and budget. This outcome supplies neither
    an upper bound on leakage nor a privacy guarantee for the whole system.
  \item[Undetermined.] Missing coverage, incomplete qualification or
    unavailable admissible observations prevents the required conclusion.
    The endpoint cannot contribute a privacy pass.
  \item[Invalid.] A demonstrated defect disconnects the measurement from
    the claimed estimand. For example, labels for the next token cannot
    establish resistance to reconstruction of the current token. The
    original score may still describe another clearly named quantity.
  \item[Out-of-scope.] The claim requires an observer, access to secrets,
    target or protocol different from the one evaluated. Even an accurate
    measurement in the original setting cannot establish the changed claim.
\end{description}

We record utility and remote gain separately and preserve each campaign's
aggregation rule. In the privacy, utility and gain exploration, a qualified
privacy breach or an independent utility failure can reject a configuration
whose other privacy endpoints remain undetermined. Incomplete qualification
blocks a conjunctive pass without erasing valid negative evidence. Failure
to meet an exploratory gain rule also leaves the population gain uncertain.

\paragraph{Transferring a conclusion.}
A conclusion can carry to another setting when the relevant protocol,
observer access, observations, target and measurement binding match, or
when an explicit argument justifies the change. A result for a wire
observer therefore needs further justification before it can support a
claim about a compromised host. Empirical non-detection provides no
information theoretic ceiling, and correlated rows add no independent
training seeds. The cases illustrate these limits on inference. They
provide neither an independent validation cohort nor evidence for a
universal privacy decision procedure.

\section{Four campaign case studies}
\label{sec:cases}

\subsection{The decoy-gradient experiment}
\label{sec:f1}

The decoy-gradient experiment mixes real activation rows with decoys.
The trusted client sends protected activations to the cloud and receives
its output. The client holds the private loss and returns the output
gradient to the cloud. Decoys receive exactly zero gradient because the
loss ignores them. That support pattern reveals which rows bear the loss,
so the disclosure concerns the real/decoy partition. Token identities
require a separate content attack~\cite{paper1} (\artifact{E1}).

The experiment records report exact partition recovery on 4,096 of 4,096
inspected frames \emph{per run} across nine seeds. A frequent token content
probe separately improves on a constant guess baseline by $+0.65$ to
$+1.50$ percentage points. The two results use different denominators.
In six cells that satisfy the utility rule and expose gradients, three
seeds on each of two datasets pass the privacy gate for forward
observations alone. The joint view exceeds the privacy gate of $+1.0$
percentage points set in advance (\artifact{E1}).

The result for forward observations remains valid within that view.
Extending its non-detection to the cloud's joint view exceeds the scope of
the test, where the additional observation exposes a breach. The records
also report that clipping each gradient row and adding Gaussian noise
removes the scored partition disclosure. The partition is disclosed in
0 of 1,024 inspected frames, with agreement 0.400 at a 32/80 base rate.
The associated held out cross entropy \emph{utility cost} is approximately
$0.01$ nats. This quantity measures utility rather than residual leakage.
The mitigation result covers the scored probes, with attacks accumulating
observations across training and other omitted families still untested.
It nevertheless supplies evidence against dismissing all perturbation
based mitigation.

\subsection{The boundary-replacement audit}
\label{sec:f2}

The boundary-replacement audit found that representations of
$(x_1,\ldots,x_{T-1})$ were scored against labels
$(x_2,\ldots,x_T)$.
A row representing $x_t$ was therefore evaluated against $x_{t+1}$.
The scores measure predictive information about the next token, while
the privacy claim concerned reconstruction of the represented input.
Even a high positive control score under this alignment validates only
the predictive target. The scores can be retained as next token probe
results, but their use to establish current token privacy is invalid
(\artifact{B1}).

Earlier independent audit notes report that separately trained controls
with a zero return match or outperform the trained remote model in all
six clean comparisons pairing corpus and seed. We retain this result as
\emph{source-reported}. Learned remote computation was therefore not shown
to be necessary for the observed utility endpoints. Removing the return
from an already trained model is a different intervention, and can affect
that model even when separate training without a return meets the utility
target (\artifact{B1}).

One matched checkpoint evaluation also exposes a difference between
utility metrics. NLL rises by approximately $0.034$ nats and passes its
rule, while teacher forced answer token top-1 match falls from $34.79\%$
to $18.15\%$. The drop is $16.64$ percentage points. This specific measure
of answer token matching differs from standard multiple choice task
accuracy, and the observation comes from one checkpoint rather than
replicated evaluation of the deployed protocol. The records therefore
narrow the quality and delegation claims while leaving corrected
reconstruction results unavailable (\artifact{B1}).

\subsection{The privacy, utility and gain exploration}
\label{sec:f3}

The privacy, utility and gain exploration completed 232 runs and rejected
all 19 configurations under its frozen exploratory rules. Reconciliation
of the cell table yields 13 qualified privacy breaches, five utility gate
failures and seven learning gain failures. The rejection reasons overlap
across configurations. Six configurations have undetermined privacy and
independent rejection on utility or gain, with width-64 rejected on gain
alone (\artifact{R1}).

Four of the 24 required instruments qualified. These were input byte
probes on inbound and outbound gradients in the two recorded defence
categories. Missing coverage prevents a joint scoped privacy pass, but
the applicable qualified probes still establish breaches. Independent
utility and gain failures also retain their force. The records thus
describe an empty feasible set for the finite tested family under its
recorded rules, alongside incomplete privacy coverage (\artifact{R1}).

The span-6 cell shows why rejection reasons must be read with the exact
inequality. Its mean utility cost across four seeds is $0.0999705344$
nats, with an upper 95\% Student $t$ limit of $0.1195024415$.
Both pass their respective thresholds of $0.10$ and $0.12$. The mean is
$2.95\times10^{-5}$ nats below its threshold. Gain failure and a qualified
privacy breach account for the rejection. These intervals describe only
four seeds, without adjustment for multiple comparisons. The small cohort
limits the conclusion to the tested family (\artifact{R1}).

The proposed confirmation packet called for eight seeds. It was never
frozen, admitted or executed. Its power upper bound of $0.263772$ for the
whole family is conditional on exploratory means and standard deviations,
with substantial sensitivity to the variance assumptions. That calculation
served planning. It supplies no probability that the 19 rejections are
false and no confirmatory outcome evidence (\artifact{R1}).

\subsection{The defence-redesign pilot}
\label{sec:f4}

The defence-redesign pilot completed 66 runs under a frozen exploratory
protocol. It covered seven cells with three paired seeds per cell and
panels of 16,384 windows. No cell met the eligibility rules. All seven
pass the remote gain rule and five pass utility. The one-bit and two-bit
quantisation configurations fail utility. The composite privacy verdicts record four failed cells
and three undetermined cells. These composite verdicts retain the original
qualification policy (\artifact{P4}).

Qualified probes still provide substantial negative evidence about the
tested defences. Reported punctuation property advantages for the worst
seed range from $0.79$ to $0.89$ where qualified, against acceptance at
$0.02$. Advantages for the next byte class range from $0.097$ to $0.112$.
Punctuation uses correlation advantage, while next byte structure uses
classification accuracy advantage. The next byte is an intentional target
here. In the boundary-replacement audit, the defect arose from using it
to claim resistance to current token reconstruction. The large measured
signals support rejection of these defence settings. The ratio between a
signal and the acceptance threshold does not diagnose instrument failure
or an unattainable gate (\artifact{P4}).

The calibration constants serve distinct purposes. The positive control
floor is $0.10$ for properties and $0.02$ for structure. The null band
is $0.02$ and privacy acceptance is $0.02$. A positive control floor does
not bound undetected leakage from above. The qualification records treat
marginal null excursions as evidence about reliability, with no finding
of instrument breakage. The finite negative result therefore concerns
the tested mechanisms, strengths, budgets and qualified targets. Untested
attacks, general token privacy and universal impossibility remain outside
that conclusion (\artifact{P4}).

\section{Reconstructing the exploration and pilot decisions}
\label{sec:results}

\subsection{Calibration coverage and reasons for rejection}

The reanalysis reconstructs seed contrasts for the privacy, utility and gain
exploration from archived losses, then checks the supplied cell and seed
tables. The pilot requires a separate calculation. We recompute paired
losses and candidate/null score arithmetic from 66 run-metrics records and
126 candidate score records, with the generated verification report as a
cross-check. These calculations use saved results throughout. Training,
probe fitting and scoring tensor captures lie outside this reanalysis.
Table~\ref{tab:cohorts} identifies the repetition units behind each summary.

\begin{table}[t]
\centering\small
\begin{tabularx}{\linewidth}{@{}lXX@{}}
\toprule
 & Exploration & Defence-redesign pilot \\
\midrule
Candidate configurations & \RoneCells & \PfourCells \\
Repetitions per candidate & 4 training seeds & 3 paired seeds \\
Recorded scope & Frozen exploratory rules & Frozen exploratory pilot \\
Qualified breach or failed composite & \RoneQualifiedBreaches & \PfourPrivacyFailures \\
Privacy undetermined & \RoneUndetermined & \PfourPrivacyUndetermined \\
Utility failures & \RoneUtilityFailures & \PfourUtilityFailures \\
Gain failures & \RoneGainFailures & 0 \\
Eligible candidate & None & None \\
\bottomrule
\end{tabularx}
\caption{No candidate meets all three requirements.
Exploration rejection reasons overlap. The pilot retains its recorded
composite privacy rule, so the two privacy rows represent different
estimands. Counts enumerate archived candidate configurations, with no
population frequency interval. The source records are \artifact{R1} and
\artifact{P4}.}
\label{tab:cohorts}
\end{table}

Thirteen exploration configurations have a qualified privacy breach.
The other six fail independent utility or gain criteria, so all 19 remain
rejected under the frozen finite protocol. Incomplete coverage cannot
explain these rejections away. Revising the instruments would need to
change the recorded rejection reasons before it could change that
conclusion.

The statistical conventions also differ between campaigns. The exploration
uses a two-sided 95\% seed interval, while the pilot uses a one-sided
95\% upper utility limit with $t_{0.95,2}=2.919986$ over three paired seeds.
We preserve each convention. Their distinction remains visible in the
checks and disjoint exploration reason intersections in
Appendix~\ref{app:checks}, rather than being absorbed into a common
confidence calculation.

The pilot also separates atomic measurements from the campaign's composite
privacy label. A qualified structure probe can detect a large signal even
when a missing property binding leaves the all-bindings assessment
undetermined. The atomic signal remains measurable. Retaining both results
preserves the source's aggregation policy and identifies the coverage
behind its verdict. Figure~\ref{fig:campaigns} presents the two campaigns
with their different metrics kept separate.

\begin{figure}[t]
\centering
\includegraphics[width=\linewidth]{figs/campaign-gates.pdf}
\caption{Qualified breaches remain valid despite incomplete coverage.
The left panel shows overlapping rejection reasons among 19 exploration
configurations evaluated with four seeds each. The right panel gives the
recorded gate outcomes for seven pilot cells with three paired seeds.
The counts describe complete campaigns and carry no population sampling
interval. The source records are
\artifact{R1} and \artifact{P4}.}
\label{fig:campaigns}
\end{figure}

\subsection{The narrow utility margin still passes}

The exploration utility rule requires a mean difference between candidate
and baseline of at most $0.10$ nats and an upper 95\% Student-$t$ limit
of at most $0.12$ nats. Four seeds give the descriptive interval
\begin{equation}
 \bar u\;\pm\;t_{0.975,3}\,\frac{s_u}{\sqrt{4}},
 \qquad t_{0.975,3}\approx3.182446.
 \label{eq:interval}
\end{equation}
The fixed evaluation data define its scope. These intervals are
unadjusted across configurations and provide no independent corpus
replication. Figure~\ref{fig:utility} separates the two thresholds.

The span-6 configuration has mean $0.0999705344$ and upper limit
$0.1195024415$. Both recorded utility criteria are satisfied. Its mean
passes by $2.9466\times10^{-5}$ nats and its upper limit by
$4.9756\times10^{-4}$ nats. The saved rejection instead combines gain
failure with a qualified privacy breach. Depth-4 has a lower point estimate
than the base configuration. This comparison and its limited seed count
cannot establish a general depth law (\artifact{R1}).

\begin{figure}[t]
\centering
\includegraphics[width=\linewidth]{figs/r1-utility-boundary.pdf}
\caption{Span-6 passes both recorded utility criteria.
Means and descriptive two-sided 95\% Student-$t$ intervals across four
training seeds are conditional on fixed evaluation text. The $0.10$ mean
threshold and $0.12$ upper-limit threshold concern different statistics.
An interval endpoint can cross the mean threshold while the utility
criteria still pass. The source is \artifact{R1}. Intervals follow
Eq.~\eqref{eq:interval}.}
\label{fig:utility}
\end{figure}
\FloatBarrier

\section{Checking the scope of privacy claims}
\label{sec:scopechecks}

\subsection{Attack success under broader access}

Observer privileges determine which attack results transfer. The following
observation distinguishes an available successful attack from a tested
attack that finds no signal.

\begin{samepage}
\begin{proposition}[Transfer under nested observations]
Fix the target distribution, success measure and resource allowance.
Suppose observer $B$ receives everything available to observer $A$, and
$B$'s allowed attack class contains every admissible attack for $A$
composed with projection to $A$'s view, within the stated resource
allowance. Then $B$ can achieve at least the success of any particular
attack available to $A$. Non-detection by a tested attack for $A$ supplies
no upper bound on the success available to $B$.
\end{proposition}
\begin{proof}
Observer $B$ projects its observations to $A$'s view and executes the same
attack. Its output distribution and success remain unchanged. The
construction shows that an attack is available. A particular implementation
or optimiser may still fail to discover it, and failure by a tested attack
supplies no upper bound.
\end{proof}
\end{samepage}

The shared-key orthogonal-wrapper experiment illustrates this distinction.
It protects against an evaluated passive wire decoder while an authorised
cloud restores the original computation. For orthogonal $Q$,
\begin{equation}
 h'=hQ,\qquad h=h'Q^{\mathsf T}.
 \label{eq:rotation}
\end{equation}
Key possession changes the observer's view. A cloud that holds or derives
$Q$ falls outside the wire observer's no-key assumption, and a compromised
host can inspect that key and invert the transform while remaining passive.
The source account and later collection manifest expressly scope the
restoration result to the wire and logged captures (\artifact{E8},
\artifact{ER}). The measured result remains valid there, while
confidentiality from the restoring host requires a different argument.

Orthogonal wrapping also preserves norms and inner products. An invertible
coordinate change leaves information recovery unresolved when a fixed
decoder fails. The result describes the tested decoder. It provides
no universal division between structure and statistics, or upper bound on
information retained in the transcript. The robustness record itself
includes stronger labelled-pair attackers and depth-dependent recovery
(\artifact{E8}).

\begin{table}[t]
\centering\small
\begin{tabular}{@{}rrrr@{}}
\toprule
Depth & Estimated $K_{50}$ & Measured bracket & $K_{50}/5120$ \\
\midrule
1 & 258.32 & 200 to 300 & 0.05045 \\
4 & 355.81 & 300 to 400 & 0.06949 \\
8 & 617.93 & 600 to 700 & 0.12069 \\
\bottomrule
\end{tabular}
\caption{Recovery crossings vary with depth at fixed width.
The 27B/$H=5120$ artifact estimates absolute 50\% token-recovery crossings
by interpolation between measured $K$ values. Uncertainty in these
conditional crossing estimates was not quantified.
The 2.39-fold depth spread at fixed width does not establish a common
$0.09H$ coefficient. The source is \artifact{E8}.}
\label{tab:k50}
\end{table}

The pooled depth-1 session-rotation score is $3.71\%$, while individual
sessions include $5.73\%$, $3.65\%$ and $2.34\%$. The distinction changes
what the score means. A pooled value describes the aggregate and cannot
serve as a worst-session confidentiality bound. The empirical crossings in Table~\ref{tab:k50}
depend on the attack. Useful partial recovery can occur below the absolute 50\%
crossing, so these laboratory curves alone cannot determine a safe key
lifetime (\artifact{E8}).

\subsection{Comparing rotation cadences}

The rotation deployment study separates a three-session Mistral-7B
comparison from a later eight-run Mistral-NeMo-12B per-epoch matrix.
Fixed and per-session rotated keys have negligible measured throughput
differences in the first comparison. In the later matrix, identity is
correct in 20 of 24 sessions and four of eight runs are admitted with known
limitations. The N=256 setting advances no epochs. Its short request horizon makes
it a no-rotation control. Wire attacks on rotated epochs
record zero top-1 recovery within the tested decoder and capture scope
(\artifact{ER}).

The faster rotation carries a measured cost. For N=16, RTT=0 and session~1,
the identity record contains $2.8$ tokens/s for the ratchet and $7.6$
tokens/s for the static arm, a $63.16\%$ throughput reduction at the
recorded precision. The ledger reports about $5.5$ to $5.6$ tokens/s for
the evidence-triggered budget arm. We retain run and session
timings separately. Their extraction adds no new paired runtime estimate,
and models, cadences and RTT twins remain separate observations rather
than independent pooled replications (\artifact{ER}).

The historical records require the same care. An earlier narrative marks
the deployed matrix invalid, whereas a later manifest admits supporting
results after further analysis. The later disposition governs our reading.
Its identity, runtime and observer limitations remain attached to those
results, preserving both the revised assessment and the history from
which it arose.

\subsection{Reading mutual information estimates}

The MI-estimator analysis records near-zero estimates on some defended
boundaries (\artifact{MI}). At the population level, the variational
objective of Donsker and Varadhan for a critic $T$ satisfies
\begin{equation}
 L_T=\mathbb E_{P_{XZ}}[T(X,Z)]-
 \log\mathbb E_{P_XP_Z}[e^{T(X,Z)}]\le \operatorname{MI}(X,Z).
 \label{eq:dv}
\end{equation}
Here $\operatorname{MI}(X,Z)$ denotes mutual information. The bound points
in one direction. Finite-sample estimation can be difficult
\cite{song2020,mcallester2020}, and the population inequality itself
cannot supply an upper privacy budget.

Let $X$ be uniform on $m$ symbols and let $Z=X$. A constant critic $T=0$
has $L_T=0$, while $\operatorname{MI}(X,Z)=\log m$ nats. With $m=256$,
the channel reveals eight bits despite the zero critic readout. The example
is purely analytic. It demonstrates the inequality's direction, without
constituting a synthetic benchmark or estimating the defended system's
actual leakage. Even a valid zero lower bound leaves a small upper bound
unestablished. A second weak lower-bound critic supplies the same limited
inference, while an empirical plug-in estimate also needs statistical
justification before it can be called a certified population bound. The
source's own limitations state that neither estimator is an upper bound
(\artifact{MI}).

We therefore retain the estimator-specific observations while withholding
a small-leakage certificate. The inequalities alone leave actual mutual
information unresolved. They also cannot establish high leakage in every
defended cell. Capacity claims need comparable care.
Similar results across finitely tested architectures and schedules show
empirical saturation. An information-theoretic ceiling requires a bound
over all remote computations.

\section{Relationship to existing evidence and methods}
\label{sec:related}

The decoy-gradient experiment and the gradient gain experiment provide
evidence about observation coverage and the units used to score
recovery~\cite{paper1,gradients2026}. In the controlled GPT-2 comparison,
exact document recovery rises from $13.71\%$ to $37.77\%$. Token recovery
is a different measure. We reuse the cited results as evidence for the
audit, with their original experimental attribution and scoring units.

The LLM-Attack Suite already specifies observation contracts, controls,
saved predictions, disclosures of untested attacks and checks on measurement
conformance~\cite{suite2026}. Those methods supply important context.
Our contribution is the retrospective reconciliation of claims across
campaigns and levels of inference. The contract terminology organises
reporting, while the five outcomes distinguish what each set of records
supports. We have neither measured their diagnostic accuracy independently
nor compared a new attack suite against the existing one for superior
coverage.

The split federated learning study by Jin et al.~\cite{privpair2026} considers a
server that observes activations and gradients. It argues that perturbation
alone provides insufficient protection in that setting and proposes
client adapters that obfuscate and recover representations. Its scope
provides a close comparison for the training setting without establishing
a universal division between value and structure. The image reconstruction
study by Smith et al.~\cite{mia2026} examines how attack and evaluator choices affect
measured leakage. It motivates careful evaluation, but leaves the qualified
breaches in the privacy, utility and gain exploration and defence-redesign
pilot intact. Limitations of MI estimators~\cite{mcallester2020,song2020}
also constrain the interpretation of critic scores, whatever defence
produced the observations.

\section{Limitations and implications for further evaluation}
\label{sec:limits}

\paragraph{Selection and dependence.}
We selected the cases from the original synthesis after their outcomes
were known. They share code, corpora and research decisions. This
dependence and the retrospective selection prevent the case counts from
estimating how often evaluation failures occur across the field.
Correcting these known examples also leaves the audit's sensitivity and
specificity on unseen studies unmeasured. A prospective independent cohort
would be needed to assess that performance.

\paragraph{Evidence available for recomputation.}
The secondary analysis reconstructs selected statistics and decisions
from supplied artifacts. The decoy-gradient experiment and parts of the
boundary-replacement audit remain reported through cited experiment
records or prior audit notes. Their attackers were not retrained, and
the new script does not recover every frame. Source hashes identify the
files consumed by the analysis, while uninspected collection history
retains its original uncertainty. Quantitative headlines exclude records
whose primary derivation could not be established.

\paragraph{Qualification and uncertainty.}
Qualification covers a specified measurement binding. It provides no
assurance against an unbounded adversary. The exploration uses four seeds
and the pilot three paired seeds, with both summaries conditional on fixed
evaluation panels and limited in precision. No new confirmation was
executed. The separate power calculation remains conditional planning.
Policy thresholds, positive control floors and the smallest reliably
detectable injected leak have different meanings and may use different
units, even when numerical constants coincide.

\paragraph{Security and utility.}
Every non-detection belongs to a finite attack and budget. A low pooled
recovery score, an invertible wrapper whose key is shared with the cloud,
or a near zero lower bound cannot establish privacy from a compromised
host. Utility also depends on the endpoint. Restoration on the tested
prompts leaves broader quality questions open, and NLL over short contexts
does not capture all of them. Noise and quantisation must be assessed at
specified strengths and utility budgets before their effectiveness can be
judged.

The audit keeps experimental outcomes and their interpretation together.
An independent validation could test whether these distinctions detect
known evaluation faults while preserving correctly detected breaches.
Such a study would need to freeze its data, budgets, attack adaptation
and primary comparison before its outcomes were observed. That validation
remains separate from the results reported here.

\section{Conclusion}

The programme records an omitted gradient channel, a mismatch between
target and claim, incomplete qualification and valid negative defence
results. The audit preserves those distinctions. The privacy, utility and
gain exploration retains 13 qualified breaches despite its 20 unqualified
instruments. The defence-redesign pilot's seven ineligible cells also give
no reason to call its gate broken. Further checks limit what the
broader security claims can support. Non-detection for a wire observer
cannot establish protection from a compromised host, and a near zero MI
lower bound cannot certify a small leakage budget. The reproducible
contribution is an account of which conclusions follow from the available
records, with the source and limits of each conclusion made explicit.

\paragraph{Data and code availability.}
The revision directory contains code for the secondary analysis, derived
tables, scripts for vector figures, source hashes and build instructions.
The relative paths in Appendix~\ref{app:artifacts} locate primary inputs
in the accompanying experiment collection. They describe the supplied
package without claiming a public repository URL or access to every
original execution environment. The packaged arithmetic can be rerun
without new training or attacks. Repeating the historical GPU experiments
is a separate reproduction task~\cite{artifacts2026}.

\appendix
\section{Sources and reproducibility}
\label{app:artifacts}

Paths below are relative to the repository root, the parent of
\path{revision-2/}. The complete list of consumed files, byte lengths
and SHA-256 values appears in \path{revision-2/data/source_manifest.json}.
The flat index is \path{revision-2/data/source-index.csv}. The manifest
binds 118 inputs to the secondary-analysis script. These records identify
the evidence consumed here, alongside the original campaigns' collection
and admission records, which retain their separate provenance role.

\begin{description}
\item[\hypertarget{artifact-E1}{Decoy-gradient experiment.}]
\path{evidence/decoy-gradient/paper1-main.tex}, especially the abstract,
mechanism section and mitigation results. The public preprint reports the
frame counts, content-probe effect and utility cost \cite{paper1}.
Those measurements are source-reported here, without a fresh count of
all frames.

\item[\hypertarget{artifact-B1}{Boundary-replacement audit.}]
\path{evidence/boundary-replacement/independent_evidence_review_2026-09-09/VERDICTS.md}
and its accompanying recomputed-evidence JSON. These records trace the
target-alignment defect and utility/control comparisons to their original
runner and outputs. Our analysis hashes those sources. Corrected attacker
fitting and the prior audit's full set of checks remain outside this
secondary analysis.

\item[\hypertarget{artifact-R1}{Privacy, utility and gain exploration.}]
\path{reproduction/trilemma-campaign/evidence/r1-exploration-20260917/analysis/exploratory_analysis.json},
\path{reproduction/trilemma-campaign/evidence/r1-qualification-20260916/attack-spec-frozen.json}, and
\path{evidence/trilemma-campaign/r1_cell_results.csv} and
\path{evidence/trilemma-campaign/r1_seed_results.csv}.
The script reconstructs 76 contrasts across configurations and seeds,
per-binding privacy rules, utility/gain decisions and consistency with
the saved tables.
\path{reproduction/trilemma-campaign/evidence/r1-exploration-20260917/confirmation_power.json}
provides conditional planning evidence for a confirmation that was never
executed.

\item[\hypertarget{artifact-P4}{Defence-redesign pilot.}]
\path{reproduction/defense-redesign-pilot/outputs/runs/pilot-20260922-a/verification.json},
the adjacent report and run-metrics records, and
\path{evidence/defense-redesign-pilot/PILOT_DECISION.md}.
The script recomputes paired utility/gain statistics, candidate advantages
and recorded qualification/composite decisions from saved candidate arrays
and qualification scores. Tensor captures remain outside this rescoring.

\item[\hypertarget{artifact-E8}{Orthogonal-wrapper experiment.}]
\path{reproduction/structural-e8/paper-data/collected/h100-2/issues/issue_e8_27b_20260807/e8_robustness_27b_bf16_densek_20260807T230529Z.json}
and \path{evidence/structural-e8/EXPERIMENT_STORY.md}.
The analysis extracts recorded baselines, crossings at each depth and
pooled/session values. The known-pair curve comes from the original
experiment and is retained at that evidential level.

\item[\hypertarget{artifact-ER}{Rotation deployment study.}]
\path{reproduction/structural-e8/paper-data/collected/toula/er_20260811/COLLECTION_MANIFEST.json},
its admitted summaries and session identity records, and
\path{evidence/structural-e8/EVIDENCE_EXECUTION_LEDGER.md}.
The script extracts eight runs and 24 sessions. For N=16, RTT=0 and
session~1, the identity record reports $2.8$ and $7.6$ tokens/s. Their
ratio gives a $63.16\%$ throughput reduction at the source's recorded
precision, with no additional hardware-measurement accuracy guarantee.

\item[\hypertarget{artifact-MI}{MI-estimator analysis.}]
\path{evidence/privacy-contract/privacy-contract-draft.tex} and
\path{evidence/privacy-contract/KEY_NUMBERS.md}, with
\path{evidence/structural-e8/LATENT_CAMPAIGN_REPORT_V7_V9.md} for the
finite capacity ablations. The near-zero scores are source-reported.
Equation~\eqref{eq:dv} and the constant-critic counterexample provide
analytic arguments, separate from the recorded model measurements.
\end{description}

\section{Checking the derived records}
\label{app:checks}

From \path{revision-2/}, the analysis runs with
\begin{quote}\small\ttfamily
python3 analysis/recompute.py
\end{quote}
The script uses the Python standard library and writes derived records only
under \path{revision-2/data/}, without loading checkpoints, tensor
captures or secret keys. Its stored log reports 1,522 arithmetic and
disposition checks on this input set. The checks reconstruct seed
contrasts, confidence endpoints, per-binding rules and composite outcomes,
then compare them with the previously generated tables. They test internal
consistency. Independent experimental replication and accuracy on unseen
studies require separate evidence.

\begin{table}[h]
\centering\small
\begin{tabular}{@{}lr@{}}
\toprule
Rejection reason combination & Configurations \\
\midrule
Qualified privacy breach only & 10 \\
Gain failure and qualified privacy breach & 3 \\
Utility and gain failure & 3 \\
Utility failure only & 2 \\
Gain failure only & 1 \\
\midrule
Total & 19 \\
\bottomrule
\end{tabular}
\caption{The reason intersections account for the full exploration.
These disjoint categories explain the overlapping marginal counts in
Figure~\ref{fig:campaigns}. The six configurations without a qualified
privacy breach have independent utility or gain rejection reasons.
The source is \artifact{R1}. Derived records appear in
\texttt{data/audit\_results.json}.}
\label{tab:intersections}
\end{table}

The derived JSON retains uncertainty conventions, denominators, source
levels and unresolved historical limitations. Before writing vector PDFs,
the figure script checks the generated CSVs against the original cell/seed
tables and pilot report. The revision README gives the regeneration
commands and dependencies. A separate preservation manifest checks the
existing \path{paper/} files independently of this revision.

\FloatBarrier
\bibliographystyle{plain}
\bibliography{refs}
\end{document}
